\scanpage{81}
\mspage{87}

Voici un \uuline{contre-exemple} que je donne \inserted{d'abord} dans le contexte des espaces topologiques, au lieu de $\mathrm{Cat}$. Je vais construire une fibration loc$^{\mathrm t}$ triviale et non triviale

% slanted
\margin{Voir un contre-exemple encore plus simple \unsure{pages|page} 93, 94. (Mais qui a l'inconvénient de ne pas être à fibres \unsure{connexes}.)}
\diagram{\begin{tikzcd}[ampersand replacement=\&,column sep=small]
X\arrow[d,"f"'] \&[-8pt] {\supset S^1}\\
S=S^1 \&
\end{tikzcd}}
de fibre $S^1$, telle que $f$ soit une $\mathbb D$-équivalence $\forall\,\mathbb D=\mathrm{Der}(k\text{-}\mathrm{Mod})$, $k$ étant un anneau où 2 est inversible.

% slanted
\margin{\uline{NB} Exemple beaucoup \deleted{\unsure{en|et}} plus simple encore : on prend $X=S^1$, et $X\to S$ le revêtement à 2 feuillets. Si on prend le rev. à \unsure{$n$|$-$} feuillets c'est \unsure{encore} OK si \unsure{$n$|$2$} est \unsure{inv.}}

\deleted{\unsure{Donc|Dan}} Cela signifie que
\[\deleted{\unsure{H|\mathbb Z}}H_{\unsure{\bullet|*}}(X,\deleted{\unsure{k}}\mathbb Z)\longrightarrow\deleted{\unsure{H|\mathbb Z}}H_{\bullet}(S,\deleted{\unsure{k}}\mathbb Z)\]
est un iso mod 2-torsion (noyau et conoyau sont annulés par une puissance de 2 -- en fait par 2 seul), ou encore $H_{\bullet}(X,k)\xrightarrow{\sim}H_{\bullet}(S,k)$ si 2 inv. dans $k$. Mais la fibre $S^1$ n'est pas $\mathbb D$-contractile (sauf si $k=0$ !)

\deleted{L'idée est de choisir la fibration telle} La fibration est donnée, à isom. \inserted{unique} près, en se fixant un point $s\in S$, et en se donnant \deleted{\unsure{alors|alor}} un autom. \inserted{$u$} de la fibre $X_s$ \deleted{\unsure{qui soit}} (``monodromie'') (la classe d'isom. du fibré ne dépend que de la classe d'isotopie \unsure{$u$|$\sim$}). Je prendrai pour $u$ la symétrie (renversant l'orientation) $z\longrightarrow z^{-1}$ du cercle \unsure{standart|standard} $\{z\in\mathbb C\mid |z|=1\}$.

Cette fibration ``se réalise'' \deleted{\unsure{par}} \deleted{\unsure{elle|de}} en un ordonné \inserted{$X\to S$}, $S^1$ étant représenté (\unsure{seule base|seule façon} \unsure{connue} sur les fibres) par
\diagram{\begin{tikzcd}[ampersand replacement=\&,column sep=small,row sep=small]
\& \& a \& \&\\
{S={}} \& \unsure{o|0|a}\arrow[ur]\arrow[dr] \& \& \unsure{o'|0'|a'}\arrow[ul]\arrow[dl] \& {=S'\amalg_R S''}\\
\& \& b \& \&
\end{tikzcd}}
où $R$ est la cat. discrète ayant deux éléments $\unsure{0|a},1$. On définit $X$ \deleted{\unsure{en}} comme $X'\amalg_Y X''$, où $X'=X\mid S'$, \deleted{\unsure{$X\mid S$}}, $X''=X\mid S''$, $Y=X\mid R$, \deleted{\unsure{avec $X=$|que $X=$}} $Y=X\mid R$.

% slanted
\margin{NB $S',S''$, \inserted{$R=S'\cap S''$} fermés dans $S$, $X',X''$, \inserted{$Y=X'\cap X''$} fermés dans $X$.}
% archive sheet number: 80
